% Requires \usepackage{xcolor} in the document preamble.
\begingroup
\color{red}

\paragraph{25-task Long Horizon.}
We evaluate an ordered sequence of 25 household tasks in a 10-cabinet
RoboCasa365 kitchen containing more than 60 objects and fixtures. Three
cabinets are locked, and objects on the counters provide distractors and
visual references. An initial search for a teapot in the final cabinet
encourages exploration; subsequent tasks reuse earlier observations and
interaction outcomes without resetting memory or the environment. The
sequence spans four task families:

\begin{itemize}
\item \textbf{Finding (9 tasks).}
Locate objects using remembered locations, neighboring items, or
accessibility. Examples include distinguishing a stored mug from one on
a counter, identifying ketchup by its surrounding condiments or fruit,
and finding the bottle with clear space for grasping.

\item \textbf{Matching (9 tasks).}
Distinguish three visual variants within each of seven object categories,
using either a specified identity or a countertop reference. Similar
colors and category labels make branding and appearance important.
Targets are also placed in unexpected locations, such as Corn Flakes in
the fruit cabinet, to discourage shortcuts based on cabinet categories.

\item \textbf{Manipulation (3 tasks).}
Select where bread, a mug, or a bowl belongs using remembered cabinet
contents. For example, bread belongs with the pantry's dry goods even
though no cabinet already contains bread. The robot selects a destination
without transporting the object.

\item \textbf{Memory queries (4 tasks).}
Recall previously observed contents and spatial relationships, such as
what was stored with the mayonnaise or where the croissant near the
pineapple was located. These questions can be answered from existing
memory without further exploration.
\end{itemize}

\endgroup
